% ============================================================================
\section*{Appendix: Proofs of Irreducibility and Reducibility}
\label{sec:appendix-proofs}
\addcontentsline{toc}{section}{Appendix: Proofs of Irreducibility and Reducibility}

The Irreducibility Criterion of \Cref{sec:irreducibility} admits or rejects a candidate element by asking whether its cost can be expressed as a same-layer dataflow composition of constituent costs. This appendix exercises the criterion on three cases: it proves irreducibility for the MAC Unit and for Attention, and proves reducibility for the Systolic Array. The three together show the criterion is neither trivially permissive (everything admitted) nor trivially restrictive (nothing admitted).

\paragraph{Proof of irreducibility: MAC Unit (Ma).} At the Hardware layer, let $\text{Ma} = g(\text{Adder}, \text{Multiplier})$. A dataflow-graph composition yields:
\begin{align}
\label{eq:mac-proof}
&f\bigl(\text{Cost}_L(\text{Adder}),\, \text{Cost}_L(\text{Multiplier})\bigr) \notag \\
&\quad = N_{\text{cores}} \cdot f_{\text{clk}} \cdot 2 \;\approx\; 67\;\text{TFLOP/s}.
\end{align}
But $\text{Cost}_L(\text{Ma}) = \cvalint{\HhPeakTFLOPS}\;\text{TFLOP/s}$ (FP16 tensor cores). The $14.8\times$ gap arises from the tensor-core pipeline's $4{\times}4{\times}4$ warp-level matrix multiply, a microarchitectural optimization invisible at the arithmetic level. Since $\cvalint{\HhPeakTFLOPS} \neq 67$ under any steady-state dataflow composition of adders and multipliers, the MAC is irreducible. \qed

\paragraph{Proof of irreducibility: Attention (At).} At the Algorithm layer, decompose attention into two Dense Dots and a Softmax: $\text{At} = g(\text{Dd}_1, \text{Sm}, \text{Dd}_2)$. Each component's cost is expressible independently: $\text{Cost}_L(\text{Dd}_i) = O(n \cdot d^2)$ FLOPs, $\text{Cost}_L(\text{Sm}) = O(n)$. But their composition creates an $O(n^2 d)$ intermediate matrix $S = QK^T$ whose memory cost is a property of the attention pattern, not of the individual matmuls:
\begin{equation}
\label{eq:attn-proof}
\text{Cost}_L(\text{At})_{\text{mem}} = O(n^2) \neq O(n d^2) + O(n) + O(n d^2)
\end{equation}
No dataflow graph over the individual costs recovers the quadratic memory term. Attention is irreducible. \qed

\paragraph{Proof of reducibility: Systolic Array.} A Systolic Array \citep{kung1982systolic} is a regular grid of $N \times N$ MAC units connected by nearest-neighbor links. Its cost decomposes as:
\begin{align}
\label{eq:systolic-proof}
\text{Throughput}(\text{SA}_{N \times N}) &= N^2 \times \text{Throughput}(\text{Ma}) \\
\text{Latency}(\text{SA}_{N \times N}) &= (2N - 1) \times \text{Latency}(\text{Ma}) \notag \\
\text{Energy}(\text{SA}_{N \times N}) &= N^2 \times \text{Energy}(\text{Ma}) + E_{\text{link}} \notag
\end{align}
where $E_{\text{link}}$ is the nearest-neighbor interconnect energy, itself a same-layer quantity. All three cost metrics are steady-state dataflow compositions of same-layer elements, assuming 100\% utilization and ignoring dataflow-specific pipeline fill/drain effects. Under those assumptions, the Systolic Array is reducible and \Cref{eq:irreducibility} is satisfied. It is rejected from the taxonomy. \qed

\paragraph{Caveat on parameterization.} The Systolic Array result is parameterized by what one admits as a same-layer element. If nearest-neighbor interconnect energy $E_{\text{link}}$ were treated as an emergent property of wire-delay variation rather than a same-layer Hardware element, the Systolic Array would reclassify as irreducible. Similarly, if dataflow-specific pipeline fill/drain effects or utilization ceilings were treated as emergent rather than bookkept at the same layer, the conclusion could flip. The criterion records a granularity choice; it does not decide it. Choosing differently is equivalent to choosing a different taxonomic resolution, and the three proofs above would then rerun against the new resolution. What the criterion does decide is that the choice is explicit, auditable, and replicable once fixed.
